\textbf{Phần 1. Trắc nghiệm lựa chọn}
\setcounterpageref{ex}{0}

\begin{ex} %[3H1V16]
Cho đoạn thẳng $AB$ dài $48\text{ cm}$. Điểm $M$ là trung điểm của $AB$, điểm $N$ là trung điểm của đoạn thẳng $AM$. Độ dài đoạn thẳng $NB$ là
\choice
{$24\text{ cm}$}
{$30\text{ cm}$}
{$40\text{ cm}$}
{\True $36\text{ cm}$}
\loigiai{
Độ dài đoạn thẳng $AM$ là: $48 : 2 = 24\text{ (cm)}$.\
Độ dài đoạn thẳng $AN$ là: $24 : 2 = 12\text{ (cm)}$.\
Độ dài đoạn thẳng $NB$ là: $48 - 12 = 36\text{ (cm)}$.
\textbf{Đáp án của câu hỏi là D}
}
\end{ex}

\begin{ex} %[3H1V16]
Trên một đường thẳng, lấy ba điểm $A, B, C$ sao cho $B$ nằm giữa $A$ và $C$. Biết $AB = 16\text{ cm}$, $BC = 16\text{ cm}$. Khẳng định nào sau đây là đúng nhất?
\choice
{$A$ là trung điểm của đoạn thẳng $BC$}
{\True $B$ là trung điểm của đoạn thẳng $AC$}
{$C$ là trung điểm của đoạn thẳng $AB$}
{$B$ chỉ là điểm nằm giữa chứ không phải trung điểm}
\loigiai{
Vì $B$ nằm giữa $A$ và $C$ đồng thời $AB = BC = 16\text{ cm}$ nên $B$ chính là trung điểm của đoạn thẳng $AC$.
\textbf{Đáp án của câu hỏi là B}
}
\end{ex}

\begin{ex} %[3H1V16]
Một sợi dây đồng dài $64\text{ cm}$ được gập đôi lại hai lần liên tiếp. Khi đó độ dài mỗi đoạn dây nhận được là
\choice
{$32\text{ cm}$}
{$8\text{ cm}$}
{\True $16\text{ cm}$}
{$12\text{ cm}$}
\loigiai{
Gập đôi lần thứ nhất, chiều dài mỗi phần là: $64 : 2 = 32\text{ (cm)}$.\
Gập đôi lần thứ hai, chiều dài mỗi phần là: $32 : 2 = 16\text{ (cm)}$.
\textbf{Đáp án của câu hỏi là C}
}
\end{ex}

\begin{ex} %[3H1V16]
Một con đường thẳng dài $100\text{ m}$. Người ta trồng cây ở cả hai đầu đường và tại chính giữa con đường (trung điểm của con đường). Sau đó tiếp tục trồng thêm 2 cây nữa ở chính giữa mỗi nửa đoạn đường đó. Hỏi khoảng cách giữa hai cây liên tiếp gần nhau nhất là bao nhiêu mét?
\choice
{\True $25\text{ m}$}
{$20\text{ m}$}
{$50\text{ m}$}
{$10\text{ m}$}
\loigiai{
Con đường được chia đều làm 4 đoạn bằng nhau. Khoảng cách giữa hai cây liên tiếp là: $100 : 4 = 25\text{ (m)}$.
\textbf{Đáp án của câu hỏi là A}
}
\end{ex}

\begin{ex} %[3H1V16]
Cho đoạn thẳng $MN = 72\text{ cm}$. Gọi $I$ là trung điểm của đoạn thẳng $MN$, $K$ là trung điểm của đoạn thẳng $IN$. Độ dài đoạn thẳng $MK$ bằng
\choice
{$36\text{ cm}$}
{$48\text{ cm}$}
{$60\text{ cm}$}
{\True $54\text{ cm}$}
\loigiai{
Ta có: $MI = IN = 72 : 2 = 36\text{ (cm)}$.\
$K$ là trung điểm của $IN$ nên $IK = 36 : 2 = 18\text{ (cm)}$.\
Độ dài $MK = MI + IK = 36 + 18 = 54\text{ (cm)}$.
\textbf{Đáp án của câu hỏi là D}
}
\end{ex}

\begin{ex} %[3H1V16]
Cho ba điểm $M, N, P$ thẳng hàng theo thứ tự đó với $MN = 20\text{ cm}, NP = 30\text{ cm}$. Gọi $A$ là trung điểm của $MN$ và $B$ là trung điểm của $NP$. Độ dài đoạn thẳng $AB$ là
\choice
{$30\text{ cm}$}
{\True $25\text{ cm}$}
{$20\text{ cm}$}
{$35\text{ cm}$}
\loigiai{
Độ dài $AN = MN : 2 = 20 : 2 = 10\text{ (cm)}$.\
Độ dài $NB = NP : 2 = 30 : 2 = 15\text{ (cm)}$.\
Độ dài đoạn thẳng $AB = AN + NB = 10 + 15 = 25\text{ (cm)}$.
\textbf{Đáp án của câu hỏi là B}
}
\end{ex}

\begin{ex} %[3H1V16]
Một thanh gỗ dài $90\text{ cm}$. Bác thợ mộc cưa thanh gỗ đó tại trung điểm của nó, sau đó lấy một nửa cưa tiếp tại trung điểm để được hai mẩu gỗ nhỏ hơn. Độ dài của một mẩu gỗ nhỏ nhất là
\choice
{$45\text{ cm}$}
{$30\text{ cm}$}
{\True $22\text{ cm và } 5\text{ mm}$}
{$20\text{ cm}$}
\loigiai{
Đổi $90\text{ cm} = 900\text{ mm}$.\
Cưa lần 1 được hai thanh dài: $900 : 2 = 450\text{ mm}$.\
Cưa lần 2 được hai thanh nhỏ dài: $450 : 2 = 225\text{ mm} = 22\text{ cm } 5\text{ mm}$.
\textbf{Đáp án của câu hỏi là C}
}
\end{ex}

\begin{ex} %[3H1V16]
Cho đoạn thẳng $AB$. Điểm $C$ nằm giữa $A$ và $B$ sao cho $AC$ ngắn hơn $CB$ là $8\text{ cm}$. Biết $AB = 40\text{ cm}$. Khoảng cách từ điểm $C$ đến trung điểm $M$ của đoạn thẳng $AB$ là
\choice
{\True $4\text{ cm}$}
{$8\text{ cm}$}
{$6\text{ cm}$}
{$5\text{ cm}$}
\loigiai{
Trung điểm $M$ cách $A$ một đoạn: $AM = 40 : 2 = 20\text{ (cm)}$.\
Độ dài đoạn $AC$ là: $(40 - 8) : 2 = 16\text{ (cm)}$.\
Khoảng cách từ $C$ đến $M$ là: $CM = AM - AC = 20 - 16 = 4\text{ (cm)}$.
\textbf{Đáp án của câu hỏi là A}
}
\end{ex}

\begin{ex} %[3H1V16]
Có 4 điểm $A, B, C, D$ cùng nằm trên một đường thẳng theo thứ tự đó. Biết $AB = BC = CD = 12\text{ cm}$. Khẳng định nào sau đây là SAI?
\choice
{$B$ là trung điểm của đoạn thẳng $AC$}
{$C$ là trung điểm của đoạn thẳng $BD$}
{Khoảng cách từ $A$ đến $D$ bằng $36\text{ cm}$}
{\True $C$ là trung điểm của đoạn thẳng $AD$}
\loigiai{
Đoạn $AC = 24\text{ cm}$, đoạn $CD = 12\text{ cm}$ nên $C$ không phải là trung điểm của $AD$ (trung điểm của $AD$ cách $A$ là $18\text{ cm}$).
\textbf{Đáp án của câu hỏi là D}
}
\end{ex}

\begin{ex} %[3H1V16]
Một sợi ruy-băng dài $84\text{ cm}$. Lan muốn cắt lấy một đoạn dài $21\text{ cm}$ nhưng trong tay không có thước đo độ dài. Lan cần gấp đôi sợi dây liên tiếp bao nhiêu lần rồi cắt theo nếp gấp gần nhất?
\choice
{$1$ lần}
{\True $2$ lần}
{$3$ lần}
{$4$ lần}
\loigiai{
Sau 1 lần gập đôi, đoạn dây có độ dài: $84 : 2 = 42\text{ cm}$.\
Sau lần gập đôi thứ hai, đoạn dây có độ dài: $42 : 2 = 21\text{ cm}$.\
Vậy Lan chỉ cần gập đôi sợi dây liên tiếp 2 lần.
\textbf{Đáp án của câu hỏi là B}
}
\end{ex}

\begin{ex} %[3H1V16]
Trên tia số, điểm $A$ ứng với số $200$, điểm $B$ ứng với số $600$. Điểm $M$ là trung điểm của đoạn thẳng $AB$ sẽ ứng với số nào trên tia số?
\choice
{$300$}
{$500$}
{\True $400$}
{$350$}
\loigiai{
Khoảng cách giữa $A$ và $B$ là: $600 - 200 = 400$.\\nTrung điểm $M$ cách điểm $A$ một khoảng: $400 : 2 = 200$.\\nVậy điểm $M$ ứng với số: $200 + 200 = 400$.
\textbf{Đáp án của câu hỏi là C}
}
\end{ex}

\begin{ex} %[3H1V16]
Cho đoạn thẳng $PQ = 56\text{ cm}$. Gọi $M$ là trung điểm của $PQ$, $N$ là trung điểm của $MQ$, $K$ là trung điểm của $NQ$. Độ dài đoạn thẳng $PK$ là
\choice
{\True $49\text{ cm}$}
{$42\text{ cm}$}
{$52\text{ cm}$}
{$45\text{ cm}$}
\loigiai{
$MQ = 56 : 2 = 28\text{ cm}$.\
$NQ = 28 : 2 = 14\text{ cm}$.\
$KQ = 14 : 2 = 7\text{ cm}$.\
Độ dài đoạn thẳng $PK = PQ - KQ = 56 - 7 = 49\text{ (cm)}$.
\textbf{Đáp án của câu hỏi là A}
}
\end{ex}

\textbf{Phần 2. Trắc nghiệm đúng sai}
\setcounterpageref{ex}{0}

\begin{ex} %[3H1V16]
Xét tính đúng sai của các phát biểu về điểm ở giữa và trung điểm:
\choiceTF
{\True Nếu điểm $M$ là trung điểm của đoạn thẳng $AB$ thì $M$ chắc chắn nằm giữa hai điểm $A$ và $B$.}
{Nếu $MA = MB$ thì điểm $M$ luôn là trung điểm của đoạn thẳng $AB$.}
{\True Mỗi đoạn thẳng chỉ có duy nhất một trung điểm.}
{Một đoạn thẳng có vô số điểm nằm ở giữa hai đầu mút của nó là sai.}
\loigiai{
\begin{itemize}
    \item Ý a) Đúng: Định nghĩa trung điểm bao hàm điểm đó phải nằm giữa.
    \item Ý b) Sai: Sai vì ba điểm $A, M, B$ có thể không thẳng hàng tạo thành tam giác cân.
    \item Ý c) Đúng: Đoạn thẳng chỉ có một vị trí chính giữa cách đều hai đầu mút.
    \item Ý d) Sai: Sai vì giữa hai đầu mút của một đoạn thẳng có vô số điểm thẳng hàng nằm giữa.
\end{itemize}
\begin{itemize}
    \item a) \textbf{Đúng}.
    \item b) \textbf{Sai}.
    \item c) \textbf{Đúng}.
    \item d) \textbf{Sai}.
\end{itemize}
}
\end{ex}

\begin{ex} %[3H1V16]
Cho đoạn thẳng $AB$ dài $80\text{ cm}$. Điểm $M$ là trung điểm của $AB$, điểm $P$ là trung điểm của $AM$:
\choiceTF
{\True Độ dài đoạn thẳng $AM$ bằng $40\text{ cm}$.}
{\True Độ dài đoạn thẳng $AP$ bằng $20\text{ cm}$.}
{Điểm $M$ là trung điểm của đoạn thẳng $PB$.}
{\True Độ dài đoạn thẳng $PB$ bằng $60\text{ cm}$.}
\loigiai{
\begin{itemize}
    \item Ý a) Đúng: $80 : 2 = 40\text{ cm}$.
    \item Ý b) Đúng: $40 : 2 = 20\text{ cm}$.
    \item Ý c) Sai: $PM = 20\text{ cm}$, nhưng $MB = 40\text{ cm}$ nên $M$ không là trung điểm $PB$.
    \item Ý d) Đúng: $PB = PM + MB = 20 + 40 = 60\text{ cm}$.
\end{itemize}
\begin{itemize}
    \item a) \textbf{Đúng}.
    \item b) \textbf{Đúng}.
    \item c) \textbf{Sai}.
    \item d) \textbf{Đúng}.
\end{itemize}
}
\end{ex}

\textbf{Phần 3. Trả lời ngắn}
\setcounterpageref{ex}{0}

\begin{ex} %[3H1V16]
Cho đoạn thẳng $AB = 60\text{ cm}$. Điểm $M$ là trung điểm của $AB$, điểm $N$ là trung điểm của $MB$. Tính độ dài đoạn thẳng $AN$ theo xăng-ti-mét.
\shortans{45}
\loigiai{
\textbf{Đáp số:} $45$\\
Độ dài $AM = MB = 30\text{ cm}$. $MN = 30 : 2 = 15\text{ cm}$. $AN = AM + MN = 30 + 15 = 45\text{ (cm)}$.
}
\end{ex}

\begin{ex} %[3H1V16]
Trên đoạn thẳng $CD$ dài $72\text{ cm}$, người ta lấy điểm $E$ sao cho $CE = 18\text{ cm}$. Hỏi điểm $E$ cách trung điểm $M$ của đoạn thẳng $CD$ bao nhiêu xăng-ti-mét?
\shortans{18}
\loigiai{
\textbf{Đáp số:} $18$\\
Trung điểm $M$ cách $C$ là: $72 : 2 = 36\text{ cm}$. Khoảng cách $EM = 36 - 18 = 18\text{ cm}$.
}
\end{ex}

\begin{ex} %[3H1V16]
Một đường ống dẫn nước thẳng dài $96\text{ m}$. Người ta lắp một van khóa nước tại trung điểm của đường ống. Sau đó lắp thêm một đồng hồ đo áp lực tại trung điểm của nửa đường ống thứ nhất. Hỏi đồng hồ đo áp lực cách đầu ống nước ban đầu bao nhiêu mét?
\shortans{24}
\loigiai{
\textbf{Đáp số:} $24$\\
Nửa ống nước dài $96 : 2 = 48\text{ m}$. Đồng hồ cách đầu ống: $48 : 2 = 24\text{ m}$.
}
\end{ex}

\begin{ex} %[3H1V16]
Cho ba điểm thẳng hàng $A, B, C$ với $B$ nằm giữa $A$ và $C$. Biết $AB = 24\text{ cm}$ và $BC$ gấp đôi $AB$. Gọi $M$ là trung điểm của $BC$. Tính độ dài đoạn thẳng $AM$ theo xăng-ti-mét.
\shortans{48}
\loigiai{
\textbf{Đáp số:} $48$\\
$BC = 24 \times 2 = 48\text{ cm}$. $BM = 48 : 2 = 24\text{ cm}$. $AM = AB + BM = 24 + 24 = 48\text{ cm}$.
}
\end{ex}

\begin{ex} %[3H1V16]
Người ta chôn 9 cây cọc thẳng hàng cách đều nhau, khoảng cách giữa cọc thứ nhất và cọc thứ chín là $64\text{ m}$. Cọc thứ năm là cọc ở chính giữa. Hỏi khoảng cách từ cọc thứ nhất đến cọc thứ năm là bao nhiêu mét?
\shortans{32}
\loigiai{
\textbf{Đáp số:} $32$\\
Cọc thứ 5 là trung điểm của đoạn thẳng nối cọc 1 và cọc 9. Khoảng cách đó là: $64 : 2 = 32\text{ m}$.
}
\end{ex}

\begin{ex} %[3H1V16]
Một sợi dây thép uốn thành hình vuông có chu vi $80\text{ cm}$. Người ta duỗi thẳng sợi dây rồi cắt tại trung điểm của nó. Tính độ dài mỗi đoạn dây nhận được theo xăng-ti-mét.
\shortans{40}
\loigiai{
\textbf{Đáp số:} $40$\\
Độ dài sợi dây ban đầu là $80\text{ cm}$. Cắt tại trung điểm thì mỗi đoạn dài: $80 : 2 = 40\text{ cm}$.
}
\end{ex}
